\documentclass[12pt]{article}
\usepackage[utf8]{inputenc}
\usepackage[T1]{fontenc}
\usepackage{hyperref}
\title{Khwedam Language Resources for K-12 Curriculum Development}
\author{Compiled by Claude Sonnet 5.5 (claude-sonnet-5-5) for Andr\'as Kornai}
\date{}
\begin{document}
\maketitle

\section{Introduction}
Khwedam (Khwe, Kxoe; ISO 639-3 xuu; Glottocode kxoe1243) is a Khoe-Kwadi language of Namibia (Zambezi/Kavango area), Botswana and Angola. Speaker estimates in search results range from about 3,700 in Namibia to about 8,000 overall. The search results were thin, and I found no dictionary, standard orthography document or Bible translation for Khwedam. Dictionary and orthography hits for Khoekhoegowab concern Nama/Damara, a different language.

\section{Primary Educational Resources}
Search summaries report an adult literacy programme in Namibia with Khwe instructors, a WIMSA-supported first-language primer, and a second primer launched later. Plans were mentioned to use Khwedam as a medium of instruction in pre-primary schools. A 2000s report says English and/or Afrikaans are the media of instruction in the area's schools, with one primary school where a teacher explained subjects in Khwedam. I did not locate the primers online.

\section{Linguistic and Reference Materials}
\begin{itemize}
\item Endangered Languages Project entry for Khwe. \url{https://endangeredlanguages.com/node/95985}
\item T. G\"uldemann (1998), NIED report (Namibian education context). \url{https://www.iaaw.hu-berlin.de/de/afrika/linguistik/mitarbeiter/prof-dr-phil-tom-gueldemann/dokumente/gueldemann-1998a-nied-report}
\item Tsammalex language entry (xuu). \url{https://tsammalex.clld.org/languages/xuu-b}
\end{itemize}
A source says Khwedam was first written by Oswin K\"ohler (University of Cologne) as early as 1960.

\section{Religious and Community Resources}
No scripture or church-based materials were found. WIMSA (Working Group of Indigenous Minorities in Southern Africa) is mentioned as having supported literacy materials; I did not confirm contact details.

\section{Digital and Audio Resources}
\begin{itemize}
\item The New Humanitarian article, which discusses the Khwe community and literacy efforts. \url{https://www.thenewhumanitarian.org/node/226285}
\end{itemize}
No audio or software was found.

\section{Curriculum Development Context}
Khwedam is not known to be taught in schools. A curriculum would need to build on the existing literacy primers and the K\"ohler-based writing tradition, together with Khwe community organizations and Namibia's Ministry of Education. Further searches (``Kxoedam'', ``Khwe dictionary'', ``Khwe New Testament'') may find more.

\end{document}